\section{The degree dependence in exponential small-cut decay}
\label{sec:cut-rate-obstruction}

The factor $2^{-d}$ in the improved exponential rate is necessary up to a polynomial factor in $d$.
Parallel bundles on a cycle give an explicit example.

Let $Q_{k,m}$ replace every edge of a $k$-cycle by $m$ parallel edges, where $k\ge3$ and $m\ge1$.
Mark all $k$ vertices, each of degree $d=2m$.
For a uniform even subgraph, let $Z$ count the closed marked vertices.
Put $s=2^{1-m}$ and define
\[
 T_s=\begin{pmatrix}1-s&s\\1-s&s/2\end{pmatrix},
 \qquad
 \lambda_\pm(s)=\frac{1-s/2\pm\sqrt{1+s-7s^2/4}}2.
\]

\begin{proposition}\label{prop:cut-rate-obstruction}
For these graphs,
\begin{equation}\label{eq:bundle-cycle-moment}
 \mathbb E2^{-Z}
 =\frac12+\frac12\operatorname{tr}(T_s^k)
 =\frac12+\frac12\bigl(\lambda_+(s)^k+\lambda_-(s)^k\bigr).
\end{equation}
For fixed $m$, the exponential decay rate of the excess above one half is $-\ln\lambda_+(s)$.
As $m\to\infty$,
\[
 -\ln\lambda_+(s)=\frac{s^2}{2}+O(s^3)
 =(1+o(1))2^{1-d}.
\]
Consequently, any bound of the form
\[
 \mathbb E2^{-Z}\le\frac12+C_d e^{-c_d k}
\]
with finite $C_d$ independent of $k$ must satisfy $c_d\le-\ln\lambda_+(2^{1-d/2})$ for even $d$.
\end{proposition}

\begin{proof}
Parity at each vertex forces all bundle parities to agree.
Their common parity is a uniform bit.
Conditional on either value, the edge choices within different bundles are independent and uniform among the $2^{m-1}$ choices of that parity.

In the odd sector every bundle is nonempty, so every marked vertex is open.
In the even sector, a bundle is empty with probability $s$.
Let $X_1,\ldots,X_k$ be the independent indicators of empty bundles in this sector.
Then
\[
 2^{-Z}=\prod_{i=1}^k\left(1-\frac12X_iX_{i+1}\right),
 \qquad X_{k+1}=X_1.
\]
Summing over the Bernoulli variables gives the transfer-matrix trace $\operatorname{tr}(T_s^k)$.
The entries are $T_s(a,b)=\mathbb P(X=b)(1-ab/2)$ for $a,b\in\{0,1\}$.
This proves~\eqref{eq:bundle-cycle-moment}.

The trace and determinant of $T_s$ are $1-s/2$ and $-s(1-s)/2$.
They give the displayed eigenvalues.
For $0<s<1$, one has $\lambda_+>0$ and $|\lambda_-|<\lambda_+$.
For $s=1$, the eigenvalues are $1/2$ and zero.
Thus the dominant eigenvalue gives the stated fixed-$m$ decay rate.
Expanding its explicit formula at $s=0$ gives $\lambda_+(s)=1-s^2/2+O(s^3)$.
The last assertion follows by letting $k\to\infty$ with $d$ fixed.
\end{proof}

The exact closure probability at each vertex is $p=2^{1-d}=s^2/2$.
Thus this family has total closure mass $\mu=kp$ and excess decay rate $(1+o(1))\mu$ as $d$ grows.
It supports a possible weighted formulation in terms of closure mass.
It rules out replacing the degree dependence by a substantially larger exponential rate uniformly over these multigraphs.
